% Transcription — fonds Alexandre Grothendieck, shelfmark 10, batch 2 (pages 21-40).
% Established from the facsimile archives/batches/10/batch-02.pdf (a mirror of
% https://grothendieck.umontpellier.fr/10.pdf), per the skill
% transcribe-grothendieck.
%
% Pass: Fable 5.1 (claude-fable-5-1), 2026-09-10 — first pass, unchecked against
% the pages by a human.
%
% The batch holds two runs.
%
% Pages 22-30 continue the « Notes anciennes » of batch 1 (pages 18, 20): black
% ink on yellowed paper, the same hand and the same objects — a motive M with
% a form φ of weight p, A = End(M), the adjoint u ↦ u^φ.
%
%   22    the four rules of u ↦ u^φ, the change of form ψ = φa with a^φ = a and
%         u^ψ = a^{-1}u^φa, the involution of A determined mod inner
%         automorphisms.
%   23    the form is a finer datum than the involution; the symmetric form
%         Tr v^φ u = Tr u^φ v on A, identified with φ̌ ⊗ φ; when it is positive
%         A is an « algèbre de Riemann ».
%   24-25 the cone A⁺ of the u^φ u, the closed cone Φ⁺ ⊂ Φ^s generated by the
%         φ^u; interior elements are units and non-degenerate forms; ψ = φa for
%         interior a has the same positivity; the cone Φ⁺ is the intrinsic
%         datum, −φ gives −Φ⁺; A⁺ is its own polar for Tr vw. A margin note on
%         page 25 lists what the abstract structure should carry.
%   26    Tr vw as the contraction pairing Φ(M) × Φ(M*) through Θ_φ and the
%         transpose; the admissible cone Φ⁺(M*) and the polarity of Φ⁺(M),
%         Φ⁺(M*).
%   27-28 the criterion a)-c) for two cones P, P* to be admissible and polar;
%         « Abstraitement » — A semi-simple with a trace, Φ a bimodule under
%         A ⊗ A of rank 1 on each side, its dual Φ*, the involution recovered
%         from a symmetric φ. Page 28 is a fast draft whose connective prose
%         is largely lost; two margin notes (one on (uv)^ψ = v^ψ u^ψ).
%   29-30 « φ riemannien »; the cones Φ⁺(φ) and Φ*⁺, polar to each other by
%         the theory of Hilbert algebras; the finitely many cones P_α of Φ
%         with disjoint interiors; the invariant « nombre des cônes
%         riemanniens dans Φ ».
%
% The cursive capital he uses for the space of forms reads sometimes as Φ,
% sometimes as a looped L; it is uniformised as Φ throughout (noted once, on
% page 24). The ring above Φ and P on pages 27-28 is transcribed \mathring
% and is not defined on these pages.
%
% Pages 32-40 are a typescript, numbered 1 to 9 at the head in his hand,
% « 10. Passage à la limite projective dans les préschémas en groupes et les
% préschémas à groupe d'opérateurs » — 10.0 to 10.16, breaking off
% mid-sentence at the foot of page 40. Every leaf carries corrections in his
% hand (pencil and blue ink): renumberings (10.1 → 10.2, 10.2 → 10.3, the
% references (10.1), (10.2) corrected accordingly), substituted references
% (« EGA 8.8.2 et 8.8.3 » → « (10.0) et (10.1) »), inserted words, a
% cancelled list at the head of page 34 replaced by a boxed phrase, a
% « Définition » for a « Proposition » (10.15), a footnote call (page 38).
% Insertions are \add{}, deletions \struck{}, and his margin notes
% \marginal{}. Page 32 carries a note in his hand on the relation of these
% leaves to another pagination; page 37 the instruction « intervertir »
% swapping two circled words, transcribed as typed with a \note{}. The foot
% of page 33 is covered by a pasted blank strip under which the first
% wording of 10.4 shows through; it is not transcribed. Typing instructions
% to the typist (« aligner », a marginal T) are omitted.
%
% Page 21 is a near-blank leaf with one struck line at its head, page 31 a
% blank leaf — both absent.
%
% The dating is the inventory's own, for the whole shelfmark.
\documentclass[11pt,a4paper]{article}
% Le préambule commun à toutes les transcriptions du fonds Grothendieck.
%
% Il tient en une idée : **l'incertitude se compose, elle ne se résout pas.**
% Une transcription qui lisse un mot douteux a détruit la seule chose pour
% laquelle on la faisait — un lecteur doit pouvoir savoir, à chaque endroit,
% ce qui était sur le feuillet et ce qui a été ajouté. Les six macros qui
% suivent sont donc l'appareil critique, et non de la décoration.
%
% Les mêmes commandes sont reconnues par `scripts/render.mjs`, qui produit la
% vue de lecture du site. Une macro ajoutée ici doit l'être aussi là-bas,
% faute de quoi le rendu échouera bruyamment — ce qui est voulu : un rendu
% silencieusement faux est pire qu'un rendu absent.

\ProvidesPackage{grothendieck}[2026/08/08 Transcriptions du fonds Grothendieck]

\usepackage{amsmath,amssymb,amsthm}
\usepackage{mathrsfs}
\usepackage{stmaryrd}
% Les diagrammes commutatifs sont partout dans ce fonds : c'est la langue
% même de la géométrie algébrique de Grothendieck, et une transcription qui
% les rendrait en prose ne transcrirait rien.
\usepackage{tikz-cd}
% Français par défaut — c'est la langue des feuillets — mais l'anglais est
% chargé aussi : la traduction et le résumé sont des documents anglais, et la
% typographie française (espaces avant les deux-points) y serait une faute.
% Ces documents basculent par \selectlanguage{english} en tête de corps.
\usepackage[english,french]{babel}
\usepackage{fontspec}
\usepackage{geometry}
\geometry{a4paper,margin=25mm,marginparwidth=28mm}
\usepackage{marginnote}
\usepackage{xcolor}
% ulem, not soul. `soul`'s \st and \ul are built on a character-by-character
% scan that XeLaTeX's font machinery breaks: the document fails with an
% undefined control sequence rather than degrading. ulem does the same two jobs
% and is Unicode-engine safe. `normalem` stops it hijacking \emph, which the
% transcriptions use for Grothendieck's own emphasis.
\usepackage[normalem]{ulem}
\usepackage{hyperref}
\hypersetup{colorlinks=true,linkcolor=black,urlcolor=black,pdfborder={0 0 0}}

% --- Métadonnées du lot -----------------------------------------------------
% Reprises telles quelles par le script de rendu, qui les affiche en tête de
% la vue de lecture. Elles fixent ce que le fichier prétend couvrir : sans
% elles, une transcription flotte sans point d'attache dans un fonds de seize
% mille feuillets.

\newcommand{\folder}[1]{\gdef\grothFolder{#1}}
\newcommand{\batch}[1]{\gdef\grothBatch{#1}}
\newcommand{\leaves}[2]{\gdef\grothFirst{#1}\gdef\grothLast{#2}}
\newcommand{\foldertitle}[1]{\gdef\grothFoldertitle{#1}}
\gdef\grothFolder{?}\gdef\grothBatch{?}\gdef\grothFirst{?}\gdef\grothLast{?}\gdef\grothFoldertitle{}

% --- Le résumé --------------------------------------------------------------
%
% Il ouvre la lecture modernisée, sous le titre « Résumé », et il est composé
% en petit. Ce n'est pas un ornement : c'est le seul passage du document qui
% ne soit pas des mathématiques, et un lecteur doit pouvoir le reconnaître
% avant de l'avoir lu — puis le sauter s'il connaît déjà le sujet, ou s'y
% arrêter s'il ne le connaît pas. Le filet qui le suit marque la reprise du
% corps.

\newenvironment{resume}
  {\small\setlength{\parskip}{0.45em}}
  {\par\medskip\hrule\medskip}

% --- Mots-clés ---------------------------------------------------------------
%
% En anglais, en clôture du résumé de la lecture modernisée : le vocabulaire
% moderne sous lequel chercher ce que les feuillets construisent. Le manifeste
% du site (`npm run manifest`) les extrait pour étiqueter la cote entière —
% cette ligne est la source unique des tags, il n'y a pas de fichier à côté.
\newcommand{\keywords}[1]{\par\smallskip\noindent
  {\footnotesize\textsc{Keywords} --- \textit{#1}}\par}

% --- L'appareil critique ----------------------------------------------------

% Le numéro de feuillet, en marge. C'est l'ancre qui permet de régler un
% désaccord sur une formule en regardant *un* feuillet plutôt qu'une chemise
% de six cents. Le site s'en sert aussi pour tourner les pages du fac-similé
% quand on fait défiler la transcription.
\newcommand{\leaf}[1]{%
  \marginnote{\footnotesize\color{gray!70}\textsf{#1}}%
  \label{leaf:#1}\ignorespaces}

% Illisible : on ne devine pas. Un mot inventé qui se lit comme les autres est
% le pire résultat possible.
\newcommand{\ill}{\textcolor{red!60!black}{[\,\ldots\,]}}

% Lecture douteuse : le mot est donné, et signalé comme incertain.
\newcommand{\uncertain}[1]{\uline{#1}}

% Ajout de l'éditeur — une lettre manquante, un mot sous-entendu.
\newcommand{\add}[1]{\textcolor{blue!50!black}{[#1]}}

% Biffé par Grothendieck lui-même. Ce qu'il a rayé fait partie du document :
% une rature dit souvent où la pensée a changé de direction.
\newcommand{\struck}[1]{\sout{#1}}

% Note du transcripteur, sur l'état matériel du feuillet.
\newcommand{\note}[1]{\par\smallskip{\small\color{gray!60!black}\textit{[#1]}}\par\smallskip}

% Note marginale de Grothendieck, distincte du corps du texte.
\newcommand{\marginal}[1]{\par\smallskip\noindent\hspace{1em}%
  {\small\color{gray!60!black}#1}\par\smallskip}

% --- Titre ------------------------------------------------------------------

\AtBeginDocument{%
  \begin{center}
    {\large\bfseries Fonds Alexandre Grothendieck --- cote n\textsuperscript{o}~\grothFolder}\\[2pt]
    {\itshape\grothFoldertitle}\\[4pt]
    {\small feuillets \grothFirst--\grothLast\ (lot \grothBatch)}\\[2pt]
    {\footnotesize\color{gray!60!black}Transcription établie d'après le fac-similé numérisé
     par l'Université de Montpellier.\\
     Les lectures incertaines sont soulignées, les illisibles notées [\,\ldots\,],
     les ajouts entre crochets.}
  \end{center}
  \vspace{1em}}

\endinput


\folder{10}
\batch{2}
\pages{21}{40}
\dating{[à partir de 1958]}
\watermark{Édition de démonstration}
\foldertitle{Catégories tensorielles : notes manuscrites (s.d.), tapuscrits (s.d.)}

\begin{document}

\section*{Notes anciennes, suite (pages 22 à 30)}

\note{Suite du feuillet « Notes anciennes » du lot précédent (pages 18 et
20) : même encre noire, même papier jauni, mêmes objets — un motif $M$ muni
d'une forme $\varphi$ de poids $p$, $A = \mathrm{End}(M)$, l'adjoint
$u \mapsto u^{\varphi}$. Dans tout le feuillet, « tt » abrège \emph{tout},
« élt » \emph{élément}, « ½-simple » \emph{semi-simple}, « néc. »
\emph{nécessairement}, « nb » \emph{nombre}, et « $=$ » placé devant un nom
se lit \emph{même}.}

\page{22}
Les \uncertain{résultats} \uncertain{pour} les l'involutions
$u \mapsto u^{\varphi}$ correspondant \uncertain{à une} forme symétrique
resp. alternée (suivant \uncertain{$p$} pair ou impair). On a \ill{} :
\[
(u+v)^{\varphi} = u^{\varphi} + v^{\varphi}
\]
\[
(\lambda u)^{\varphi} = \lambda u^{\varphi} \quad (\lambda \in \mathbb{Q})
\]
\[
(u^{\varphi})^{\varphi} = u
\]
\[
(uv)^{\varphi} = \struck{\ill{}}\ v^{\varphi} u^{\varphi}
\]
\note{Le membre de droite de la dernière formule est précédé d'un facteur
biffé, non lu.}

Si on remplace $\varphi$ par une forme $\psi$ du même type (paire resp.
impaire) on aura \uncertain{$\psi = \varphi a$}
\[
\psi(x,y) \struck{\ill{}} = \varphi(ax, y)
\]
et la hyp. sur $\psi$ s'exprime comme on \uncertain{l'a dit} par
$a^{\varphi} = a$. On aura de plus
\[
u^{\psi} = \struck{\ill{}}\ a^{-1} u^{\varphi} a \qquad (a \text{ étant inversible}).
\]
donc l'involution de $A = \mathrm{End}(M)$ est \emph{bien déterminée mod
automorphismes intérieurs} (indépendamment du choix des $\varphi$, —
$\psi$), — en fait si $\varphi$ est fixé, mod aut. int. par un élément
\struck{\ill{}} $\varphi$-hermitien \struck{symétrique} de $A$.

\page{23}
On notera cependant que la donnée d'une forme $\varphi$ est plus précise
que la donnée de l'involution sur $A = \mathrm{End}(M)$, ainsi si
$\mathrm{End}(M)$ est commutatif, ttes les $\varphi$ donnent la même
involution ! De même, si $\lambda \in \mathbb{Q}^{*}$, $\lambda\varphi$
donne la même involution que $\varphi$.
\note{Un mot interlinéaire au-dessus de « forme $\varphi$ », \ill{}.}

Considérons la forme symétrique
\[
\mathrm{Tr}\, v^{\varphi} u = \mathrm{Tr}\, u^{\varphi} v
\]
\struck{sur $A \simeq \mathrm{End}(M) \simeq \check M \otimes M$, par la forme}
cette forme n'est autre que $\check\varphi \otimes \varphi$ où
(\uncertain{déduite de $\varphi$ par transposée structurale par
$\varphi$}) $\check\varphi$ est la forme duale de $\varphi$ sur $\check M$,
\struck{\ill{}}
\note{Suivent deux lignes biffées et encadrées, dans lesquelles on lit
« $\mathrm{End}(M)$ » et « $\simeq M \otimes M$ ».}
définit ainsi une forme sur $\mathrm{End}(M)$ : à valeurs dans
$\mathbb{Q}$, notée \ill{}.

Supposons que cette forme soit \struck{définie} positive (elle est néc.
non dég., donc définie positive). \struck{Soit $P$ le cône des} Alors
$A = \mathrm{End}(M)$ devient une « algèbre de Riemann ».

\page{24}
Soit $A^{+}$ un cône positif, \struck{engendré} \uncertain{l'adhérence} de
l'ens. des combinaisons linéaires à coefficients $\geqslant 0$ des
$u^{\varphi} u$ ($u \in A$ ; on peut \uncertain{supposer} $u \in A^{*}$,
groupe des unités de $A$). Il correspond dans $\Phi^{s}$ au cône fermé
$\Phi^{+}$ engendré par les $\varphi^{u}$ [$\varphi^{u}(x,y) =
\varphi(ux, uy)$], i.e. le plus petit cône fermé dans $\Phi^{s}$ contenant
$\varphi$ et stable par les automorphismes de $M$.
\note{La majuscule cursive qui désigne l'espace des formes se lit tantôt
$\Phi$, tantôt comme un $L$ bouclé ; on l'a uniformisée en $\Phi$ dans
tout le feuillet. L'exposant $s$ n'est pas défini sur ces pages
(\emph{symétriques}, vraisemblablement).}

\struck{Tout élément de} De la théor. des algèbres de Riemann résulte que
tt élt $a$ intérieur à $A^{+}$ est une unité, donc tt élt intérieur
\uncertain{à} $\Phi^{+}$ est une forme non dégénérée ; \struck{\ill{}} De
plus, par un tel $a$, \struck{l'involution} (correspondant \struck{\ill{}
$\varphi$ engendre une des cône} $\varphi \in \Phi^{+}$ intérieur
\uncertain{à} $\Phi^{+}$) l'involution
\[
u \mapsto u^{\psi} = a^{-1} u^{\varphi} a
\]
\struck{satisfait} étant déduite
\struck{$u^{\psi} u = a^{-1} u^{\varphi} a u = \sum_i b_i^{-1}
(b_i^{\varphi})^{-1} u^{\varphi} b_i^{\varphi} b_i u$}
de la précédente par un automorphisme (intérieur) \ill{}, elle
\struck{\ill{}} possède
\page{25}
les mêmes propriétés de positivité, savoir
\[
\mathrm{Tr}\, u^{\psi} u \;\bigl(= \mathrm{Tr}\, (a^{-1}ua)^{\varphi}
(a^{-1}ua)\bigr) \geqslant 0
\]
égalité seulement si $u = 0$. Donc tt élt de l'intérieur du cône
\struck{de} $\Phi^{+}$ engendre (via automorphismes de $M$) le même cône
$\Phi^{+}$. \emph{C'est la donnée du cône $\Phi^{+}$ qui, intrinsèquement,
constitue une structure supplémentaire importante.}

\marginal{\uncertain{Dégager} la structure \uncertain{déduite} :
\struck{\ill{}} $A$ \uncertain{riemannienne}, $\Phi$ sur $A$
\uncertain{libre}, \ill{}, Trace sur $A$, \uncertain{autres} ; le bon
cône de $\Phi$}
\note{Note marginale écrite en diagonale, à gauche du texte, liste rapide
dont les mots sont pour moitié perdus.}

On notera que, sous cette forme abstraite, on ne peut (parmi les cônes
possibles, correspondants aux divers choix de $\varphi$) en distinguer de
meilleurs que d'autres. Ainsi, si $\varphi$ satisfait la condition de
positivité \ill{}, de même $-\varphi$, qui définit le cône $-\Phi^{+}$.

Note aussi : \struck{Pour un cône convexe fermé} si $v, w \in A^{+}$,
\struck{\ill{}} alors $\mathrm{Tr}\, vw \geqslant 0$ [et en fait, $A^{+}$
est son propre polaire dans $A^{s}$ pour la forme $\mathrm{Tr}\, vw$].
\note{« hermitiens » ajouté au-dessus de $A^{s}$.}
Or pour
\page{26}
\struck{Or, la forme $\mathrm{Tr}\, vw$ devient la forme sur
$\Phi \simeq (\check M \otimes \check M(p))$ \uncertain{déduite} par
$\check\varphi \otimes \check\varphi$ \ill{}, ou encore,}
\note{Bloc encadré et barré d'une diagonale.}
identifiant $A$ d'une part par $\Theta_{\varphi}$ à
$\Phi(M) = \check M \otimes \check M(p)$, d'autre part via
\struck{\ill{}} $u \mapsto {}^{t}u$ et $\check\varphi$ [ou via
$\Phi(M) \simeq \Phi(M^{*})$, à l'aide de l'isom.
$\tilde\varphi \colon M \to M^{*}$] à $\Phi(M^{*}) = M \otimes M(-p)$, la
forme $\mathrm{Tr}\, vw$ sur $A \times A$ devient la forme évidente de
contraction sur $\Phi(M) \times \Phi(M^{*})$. Notons d'autre part que pour
tt cône \struck{positif} \struck{\ill{}} « admissible » $\Phi^{+}(M)$ dans
$\Phi(M)$, \uncertain{correspond} un « cône admissible » $\Phi^{+}(M^{*})$
dans $\Phi(M^{*})$, en prenant les $\check\varphi^{*}$
($\varphi \in \Phi^{+}(M^{*})$).
\note{Lecture douteuse de l'argument entre parenthèses ; on attendrait
$\varphi \in \Phi^{+}(M)$.}
Ceci posé, on voit que \emph{$\Phi^{+}(M)$ et $\Phi^{+}(M^{*})$ sont
polaires l'un de l'autre.} Conséquence :

\page{27}
Si on a deux cônes convexes \struck{fermés} $P$, $P^{*}$ dans $\Phi(M)$,
$\Phi(M^{*})$, \struck{qui sont} tels que
\begin{itemize}
  \item[a)] $P$, $P^{*}$ stables par automorphismes
  \item[b)] $(P, P^{*}) \geqslant 0$
  \item[c)] Il existe $\varphi$, $\varphi^{*}$ polaires l'un de l'autre,
    « admissibles » [i.e. \struck{Pour} $\varphi$ non dég. et
    $\mathrm{Tr}\, u^{\varphi} u \geqslant 0$ pour tt $u \in A$], avec
    $\varphi \in P$, $\varphi^{*} \in P^{*}$
\end{itemize}
alors $P$, $P^{*}$ sont \uncertain{nécessairement} deux cônes
admissibles, (et \uncertain{en particulier} chacun est \uncertain{déterminé}
par l'autre, et $\mathring{P}{}^{+}$ contient les $\check\varphi$ pour
$\varphi \in \mathring{P}{}^{+}$).
\note{Le rond au-dessus de $P$, et plus loin de $\Phi$, n'est pas défini
sur ces pages ; dans la première occurrence, un exposant est biffé et
remplacé par $+$.}

\textbf{Abstraitement}

Nous avons une algèbre $A$ ½-simple de type fini, munie d'une forme
$\mathrm{tr}$ sur $\mathbb{Q}$, telle que $\mathrm{tr}(uv)$ non dégénérée,
un \struck{bimodule} $\Phi$ (à droite) sous $A \otimes A$ (\ill{} $\rho$,
l'\uncertain{isom.} \ill{} de $A$ écrite $u\varphi v$)
\struck{transportés} de rang $1$ sur $A \otimes 1$ et sur $1 \otimes A$
\struck{(isom. correspond. à $A$)}, le $\mathbb{Q}$-dual $\Phi^{*}$,
\struck{\ill{}} muni de sa structure de module sur
$A^{\circ} \otimes A^{\circ}$ \struck{$A$ par conséquent}. Pour un élément
$\varphi$ de $\Phi$, il \ill{} de dire qu'il est \uncertain{une bonne}
\ill{} \struck{\ill{}}
\note{« sur $\mathbb{Q}$ » et « (à droite) » sont ajoutés en interligne.
Les deux dernières lignes de la page, dont une biffée, ne se lisent pas.}

\page{28}
\struck{il lui correspond alors $\tilde\varphi$ un élément canonique}
\struck{\ill{}}
\note{Deux lignes biffées ; dans la première, « unique » est biffé et
remplacé par « canonique » avant que la ligne entière ne le soit.}
sur $1 \otimes A$ — sur $A \otimes 1$. L'utilisant pour écrire un isom. de
modules sur $A \otimes 1$ : $\Phi \simeq A$, et \struck{utilisant}
composant avec $\mathrm{tr}$ sur $A$ \struck{\ill{} dont $A \simeq A$}, on
\uncertain{obtient une} forme sur $\Phi$, i.e. un élt \ill{}
$\check\varphi$. Provient ? \ill{} de forme pareil. On \struck{peut}
\uncertain{utiliser} une symétrie de $\Phi$ qui \uncertain{échange} les
deux structures de $A$-modules. L'application
$\mathring{\Phi}{}^{*} \to \mathring{\Phi}{}^{*}$ commutant néc. à cette
symétrie. Pour tt élt $\varphi$ de $\mathring{\Phi}{}^{s}$ [fixe par
symétrie] \uncertain{au moyen d'une forme} \uncertain{une bijection}
$A \to \Phi$, d'où sur $A$ une \struck{involution} application
$u \mapsto u^{\varphi}$ linéaire, telle que $(u^{\varphi})^{\varphi} = u$,
[dites \ill{} déterminée mod \struck{composition} composition avec aut.
int. par un élt $a$ de $A$ tel que $a^{\varphi} = a$, si
$\varphi \sim \psi$]. \uncertain{Ceci} on fait la \ill{} par
$\check\varphi$, — donner une \ill{} $A^{\circ} = A$ (semi-simplicité !)
\ill{} ! La structure de \uncertain{donnée d'une} \ill{} la structure
$a\,u\,b^{\varphi}$, et $\varphi \sim \varphi'$ donnent
$u \mapsto u^{\varphi}$ \ill{}
\note{Page de brouillon rapide : les formules portent, la prose de liaison
est aux deux tiers perdue.}

\marginal{\ill{} (mais non : conjugaison par !)}
\note{Note marginale cerclée, en haut à gauche, ouverte par un grand
signe, $Z$ ou $\Sigma$, non lu.}

\marginal{On a ici $(uv)^{\psi} = v^{\psi} u^{\psi}$ (indépendant du
choix particulier de $\psi$)}

\page{29}
On dit que \struck{$(A, \Phi)$} \struck{$\Phi$} $\varphi$ est
\struck{\ill{}} \emph{riemannien} si $A$, muni de la trace et de
l'involution $u \mapsto u^{\varphi}$, est une algèbre riemannienne, i.e.
\uncertain{lorsque}
\[
\mathrm{Tr}\, u^{\varphi} u \geqslant 0 \qquad \forall u
\]
\struck{\ill{}}
[\textbf{NB} utilisant les identifications $A \simeq \Phi$,
$A^{\circ} \simeq \Phi^{*}$ définies par $\varphi$ et
\uncertain{$\check\varphi^{-1}$}, la forme $\mathrm{Tr}\, u^{\varphi} u$
sur $A^{\circ} \times A$ s'identifie à la forme canonique sur
$\Phi \times \Phi^{*}$, elle est donc \emph{non dégénérée}].

Tout $\varphi$ riemannien définit comme plus haut un cône convexe fermé
$\Phi^{+}(\varphi) = \Phi^{+}$, savoir le plus petit cône convexe fermé
stable par $u \otimes u$ ($u \in A$) contenant $\varphi$, et qui correspond
par $\Theta_{\varphi} \colon A \simeq \Phi$ au cône des « éléments
positifs » de l'algèbre \struck{hilb.} riemannienne $A$, engendré par les
éléments de la forme $u^{\varphi} u$. \uncertain{De même},
$A^{\circ} \to \Phi^{*}$ transforme \ill{} en \ill{} $\Phi^{*+}$ défini
en termes de \uncertain{$\check\varphi^{-1}$}. La théorie des alg.
hilbertiennes nous dit que les cônes $\Phi^{+}$ et $\Phi^{*+}$ sont
polaires l'un de l'autre, \struck{\ill{}} que leurs
\page{30}
intérieurs \uncertain{correspondent} aux formes \ill{} non dégénérées, ou
encore (étant toutes des riemanniennes \ill{}) riemanniennes, et que pour
\struck{tt} $\psi$ de l'intérieur de $\Phi^{+}(\varphi)$ le cône
$\Phi^{+}(\psi)$ \uncertain{est égal à} $\Phi^{+}(\varphi)$.
\note{Un mot interlinéaire, « pri », au-dessus de « riemanniennes » dans
la parenthèse, non résolu.}

On trouve ainsi dans $\Phi$ un certain nombre de cônes convexes fermés
$P_{\alpha}$ deux à deux \uncertain{distincts}, stables par les
$u \otimes u$ ($u \in A$) et dont les intérieurs sont deux à deux
disjoints, la réunion des intérieurs étant l'ens. des formes riemanniennes
\uncertain{relativement} à $\Phi$.
\note{« fermés » et « des intérieurs » ajoutés en interligne.}
S'il y a une $\varphi$ riemannienne, alors $-\varphi$ est riemannienne,
mais certainement pas dans le même cône (les deux cônes \ill{} opposés).
Quel est \uncertain{la} signification de l'invariant : « nb des cônes
riemanniens dans $\Phi$ » (qu'on peut aussi considérer comme un invariant
pour une alg. riemannienne \uncertain{ou} $A$, — \ill{} sur $\mathbb{R}$
\ill{}].

\section*{Tapuscrit : « 10. Passage à la limite projective dans les préschémas en groupes et les préschémas à groupe d'opérateurs » (pages 32 à 40)}

\note{Tapuscrit de neuf feuillets numérotés $1$ à $9$ en tête, de sa main,
portant les numéros 10.0 à 10.16 ; il s'interrompt au milieu d'une phrase
au bas de la page 40. Chaque feuillet est corrigé de sa main, au crayon et
à l'encre bleue : les insertions sont en \add{}, les suppressions en
\struck{}. La machine écrit « lim » avec une flèche pour la limite
inductive, un $\mathbb{T}$ gras pour la topologie, $\mathcal{O}$ pour le
faisceau structural et $\kappa$ pour le corps résiduel ; le symbole
$\psi$ des pages 38-39 est tapé.}

\page{32}
\marginal{Les \uncertain{11} pages qui suivent sont le texte correspondant
aux \uncertain{anciennes} pages 82 à 117)}
\note{Note cerclée dans la marge gauche, de sa main.}

\subsection*{10. Passage à la limite projective dans les préschémas en groupes et les préschémas à groupe d'opérateurs.}

10.0. Rappelons le résultat essentiel de (EGA IV 8.8) : Etant donnée la
situation suivante : $S_{0}$ un préschéma quasi-compact et quasi-séparé,
$I$ un ensemble préordonné filtrant croissant, $(\mathcal{A}_{i})_{i \in I}$
un système inductif de $\mathcal{O}_{S_{0}}$-Algèbres commutatives
quasi-cohérentes, $\mathcal{A} = \varinjlim \mathcal{A}_{i}$,
$S_{i} = \mathrm{Spec}\, \mathcal{A}_{i}$ pour $i \in I$, et
$S = \mathrm{Spec}\, \mathcal{A}$, alors la catégorie des $S$-préschémas de
présentation finie est déterminée à équivalence près par la donnée des
catégories des $S_{i}$-préschémas de présentation finie, des foncteurs
$\rho_{ij} \colon X_{i} \mapsto X_{i} \times_{S_{i}} S_{j}$ pour
$i \leqslant j$ \struck{entre ces catégories}, et des isomorphismes de
transitivité $\rho_{jk} \circ \rho_{ij} \xrightarrow{\sim} \rho_{ik}$.
\struck{Plus précisément} \add{Précisons.}
\note{« entre ces catégories » est cerclé et relié par un trait au mot
« foncteurs », pour être déplacé.}

Etant donné $j \in I$, et un $S_{j}$-préschéma de présentation finie
$X_{j}$, nous poserons, pour tout $i \in I$ tel que $j \leqslant i$,
$X_{i} = X_{j} \times_{S_{j}} S_{i}$, et $X = X_{j} \times_{S_{j}} S$.
Alors \add{(EGA IV 8.8.2)} :
\begin{itemize}
  \item[i)] Etant donné $j \in I$, et deux $S_{j}$-préschémas de
    présentation finie $X_{j}$ et $Y_{j}$, l'application canonique de
    $\varinjlim_{i \geqslant j} \mathrm{Hom}_{S_{i}}(X_{i}, Y_{i})$ dans
    $\mathrm{Hom}(X, Y)$ est bijective
  \item[ii)] Pour tout $S$-préschéma de présentation finie $X$, il existe
    un indice $j \in I$, \struck{et} un $S_{j}$-préschéma de présentation
    finie $X_{j}$ et un $S$-isomorphisme
    $X \simeq X_{j} \times_{S_{j}} S$.
\end{itemize}

On en conclut (EGA IV 8.8.3) que, chaque fois qu'on aura un diagramme $D$
portant sur un nombre fini d'objets et de flèches de la catégorie des
$S$-préschémas de présentation finie, on peut trouver un indice $i \in I$
et un diagramme $D_{i}$ dans la catégorie des $S_{i}$-préschémas de
présentation finie, tels que le diagramme $D$ provienne à isomorphisme
près du diagramme $D_{i}$ par le changement de base $S \to S_{i}$. On peut
même trouver $i$ et $D_{i}$ tels que tout carré cartésien de $D$ provienne
d'un carré cartésien de $D_{i}$.

\add{10.1.} De plus, un grand nombre de propriétés \struck{\ill{}}
courantes pour un morphisme, stables par changement de base, possèdent la
propriété suivante ; étant donné un indice $j \in I$ deux
$S_{j}$-préschémas de présentation finie $X_{j}$ et $Y_{j}$ et un
$S_{j}$-morphisme $u_{j} \colon X_{j} \to Y_{j}$, pour que
$u = u_{j} \times_{S_{j}} S$ ait la propriété $\mathbb{P}$, il faut et il
suffit qu'il existe un indice $i \in I$ tel que $j \leqslant i$ et que
$u_{i} = u_{j} \times_{S_{j}} S_{i}$ ait la propriété $\mathbb{P}$. Il en
est ainsi dans le cas où $\mathbb{P}$ est \struck{la} \add{l'une des}
propriété\add{s} \add{suivantes} pour un morphisme : être séparé,
surjectif, radiciel, affine, quasi-affine, fini, quasi-fini, propre,
projectif, quasi-projectif, un isomor-
\page{33}
phisme, un monomorphisme, une immersion, une immersion ouverte, une
immersion fermée, lisse, non ramifié ou étale (EGA IV 8.10.5, 11.2.6 et
17.7.6). Remarquons qu'il en est encore ainsi \struck{\ill{}} dans le cas
où $\mathbb{P}$ est la propriété d'être couvrant pour la topologie
fidèlement plate localement de présentation finie (notée fppf) ; en effet,
étant donnés deux $S$-préschémas de présentation finie $X$ et $Y$, et un
\struck{morphisme} \add{$S$-morphisme} $u \colon X \to Y$, il résulte de
(Exp. IV 6.3.1 (i)) que, pour que $u$ soit couvrant pour la topologie
fppf, il faut et il suffit qu'il existe un $S$-préschéma $Z$ et un
$S$-morphisme $v \colon Z \to Y$ fidèlement plat et de présentation finie
qui se factorise à travers $u$.

Le but de ce numéro est \struck{d'étendre} \add{de donner des variantes
de} ce genre de résultats \struck{à} \add{pour} la catégorie des
$S$-groupes de présentation finie, \struck{et à} celle des $S$-préschémas
à groupe d'opérateurs, et \struck{à} \add{pour} certaines propriétés pour
des monomorphismes de groupes (être invariant, central, à faisceau
quotient inversible, etc ...).

Les deux résultats préliminaires de ce type sont les suivants :

\textbf{10.\struck{1}\add{2}.} Dans la situation rappelée au début de
(10.0), soient \struck{\ill{}} $j \in I$, $G_{j}$ et $H_{j}$ deux
$S_{j}$-groupes de présentation finie ; posons, pour tout $i \in I$ tel
que $j \leqslant i$, $G_{i} = G_{j} \times_{S_{j}} S_{i}$,
$G = G_{j} \times_{S_{j}} S$, et définissons de même $H_{i}$ et $H$. Alors
l'application canonique de
$\varinjlim_{i \geqslant j} \mathrm{Hom}_{S_{i}\text{-gr.}}(G_{i}, H_{i})$
dans $\mathrm{Hom}_{S\text{-gr.}}(G, H)$ est bijective.

\textbf{10.\struck{2}\add{3}.} Dans la situation rappelée au début de
(10.0), soit $G$ un $S$-groupe de présentation finie ; alors il existe un
indice $j \in I$, \struck{et} un $S_{j}$-groupe de présentation finie
$G_{j}$, et un \struck{\ill{}} $S$-isomorphisme de groupes de $G$ sur
$G_{j} \times_{S_{j}} S$.

Les assertions (10.\struck{1}\add{2}) et (10.\struck{2}\add{3}) sont des
conséquences faciles de \struck{EGA IV 8.8.2 et 8.8.3} \add{(10.0) et
(10.1)} compte tenu de l'interprétation de la structure de $S$-groupe
donnée en (EGA $0_{\mathrm{III}}$ 8.2.5 et 8.2.6).
\note{Le bas de la page est recouvert d'une bande collée, sous laquelle
transparaît une première rédaction de 10.4 (une liste (i), (ii), (iii)
de propriétés d'un morphisme de préschémas en groupes) ; elle n'est pas
transcrite. La page 34 s'ouvre sur la fin de cette liste, biffée.}

\page{34}
\struck{fermé de $H$),}

\struck{(v) être un monomorphisme central (6.6),}

\struck{vi) être un monomorphisme invariant (6.6).}
\marginal{soit un monomorphisme central (resp. un monomorphisme
invariant)}
\note{Les deux lignes (v), (vi) sont barrées de traits obliques, et la
phrase encadrée dans la marge droite, reliée par un trait au lieu de
« ait la propriété $\mathbb{P}$ » biffé ci-dessous, la remplace ; le
numéro « 10.4. » est écrit dans la marge gauche.}

\add{10.4.} Soient $j \in I$, $G_{j}$ et $H_{j}$ deux $S_{j}$-groupes de
présentation finie, et $u_{j} \colon G_{j} \to H_{j}$ un morphisme de
$S_{j}$-groupes. Pour que $u = u_{j} \times_{S_{j}} S$ \struck{ait la
propriété $\mathbb{P}$}, il faut et il suffit qu'il existe $i \in I$,
$i \geqslant j$ tel que $u_{i} = u_{j} \times_{S_{j}} S_{i}$ ait la
\add{même} propriété \struck{$\mathbb{P}$}.

\struck{Les assertions (i) à (iv) résultent de (EGA 8.10.5). Les
assertions (v) et (vi)} \add{C'est} \struck{sont} une conséquence
immédiate de \struck{EGA 8.8.2 et 8.8.3} \add{(10.0) et (10.1)} compte
tenu de la caractérisation donnée en (6.7) des monomorphismes de groupes
centraux et invariants.

\add{Corollaire} \textbf{10.5.} Soit $j \in I$, et soit $G_{j}$ un
$S_{j}$-groupe de présentation finie. Pour que
$G = G_{j} \times_{S_{j}} S$ soit abélien, il faut et il suffit qu'il
existe $i \in I$, $i \geqslant j$, tel que $G_{i} = G_{j} \times_{S_{j}}
S_{i}$ soit abélien.

En effet, il revient au même de dire qu'un $S$-groupe est abélien, ou que,
considéré comme sous-préschéma en groupes de lui-même, il est central.

\textbf{Proposition 10.6.} — Dans la situation rappelée au début de
(10.0), soient $j \in I$, $G_{j}$ un $S_{j}$-groupe de présentation finie,
$G'_{j}$ un sous-préschéma en groupes de $G_{j}$ plat et de présentation
finie sur $S_{j}$. Pour que
$G_{j} \times_{S_{j}} S / G'_{j} \times_{S_{j}} S$ soit représentable
(pour la topologie fpqc), il faut et il suffit qu'il existe $i \in I$,
$i \geqslant j$, tel que
$G_{j} \times_{S_{j}} S_{i} / G'_{j} \times_{S_{j}} S_{i}$ soit
représentable.

C'est une conséquence du lemme plus général suivant :

\textbf{Lemme 10.7.} — Soit $\mathbb{T}$ la topologie fppf ou fpqc ;
soient $j \in I$, $X_{j}$ un $S_{j}$-préschéma de
\page{35}
présentation finie, et $R_{j}$ une relation d'équivalence sur $X_{j}$
définie par un $S_{j}$-monomorphisme
$v_{j} \colon R_{j} \to X_{j} \times_{S_{j}} X_{j}$ tel que le morphisme
$\mathrm{pr}_{1,j} \circ v_{j} \colon R_{j} \to X_{j}$ déduit de $v_{j}$
par composition avec la première projection
$X_{j} \times_{S_{j}} X_{j} \to X_{j}$ soit plat et de présentation finie.
Alors, pour que le faisceau quotient
$X_{j} \times_{S_{j}} S / R_{j} \times_{S_{j}} S$ pour la topologie
$\mathbb{T}$ soit représentable, il faut et il suffit qu'il existe
$i \in I$, $i \geqslant j$, le faisceau quotient
$X_{j} \times_{S_{j}} S_{i} / R_{j} \times_{S_{j}} S_{i}$ pour la
topologie $\mathbb{T}$ soit représentable.

Compte tenu des énoncés de (EGA IV 8.8.2, 8.8.3, 8.10.5 et 11.2.6)
rappelés en (10.0), ce lemme est conséquence du résultat suivant :

\textbf{Lemme 10.8.} — Soit $\mathbb{T}$ la topologie fppf ou fpqc ;
soient \struck{\ill{}} $X$ un $S$-préschéma de présentation finie (resp.
localement de présentation finie), $R$ une relation d'équivalence sur $X$
définie par un monomorphisme $v \colon R \to X \times_{S} X$ tel que
$\mathrm{pr}_{1} \circ v$ soit plat et de présentation finie (resp. plat
et localement de présentation finie). Alors les conditions suivantes sont
équivalentes :
\begin{itemize}
  \item[(i)] Le faisceau quotient $X/R$ pour la topologie $\mathbb{T}$ est
    représentable.
  \item[(ii)] il existe un $S$-préschéma de présentation finie (resp.
    localement de présentation finie) $Y$ et un morphisme fidèlement plat
    $p \colon X \to Y$ tel que le diagramme
\end{itemize}

\begin{tikzcd}
R \arrow[r, "\mathrm{pr}_{1} \circ v"] \arrow[d, "\mathrm{pr}_{2} \circ v"'] & X \arrow[d, "p"] \\
X \arrow[r, "p"'] & X
\end{tikzcd}

\note{Le diagramme est étiqueté (D) sur la page ; son coin inférieur
droit est tapé $X$, coquille pour $Y$, laissée telle quelle.}
(où $\mathrm{pr}_{1}$ et $\mathrm{pr}_{2}$ sont les projections
$X \times_{S} X \rightrightarrows X$) soit cartésien.

Notons d'abord que d'après (Exp. IV 3.3.2 et 4.4.3), pour que le faisceau
$X/R$ pour la topologie $\mathbb{T}$ soit représentable par $Y$, il faut
et il suffit que le diagramme (D) soit cartésien et que $p$ soit couvrant
pour la topologie $\mathbb{T}$.

Montrons que i) entraîne ii). L'hypothèse (i) implique que le diagramme
(D) est cartésien, donc que $\mathrm{pr}_{1} \circ v$ se déduit de $p$ par
changement de base par $p$, et que $p$ est
\page{36}
est couvrant pour la topologie fpqc. Puisque $\mathrm{pr}_{1} \circ v$ est
fidèlement plat et de présentation finie, (resp. fidèlement plat et
localement de présentation finie), il en est de même de $p$ (EGA IV
2.7.1), et comme $X$ est de présentation finie (resp. localement de
présentation finie) sur $S$, il en est de même de $Y$ (EGA IV 17.7.4).
\note{« est » est répété au passage de la page 35 à la page 36 ; tel
quel.}

Montrons que ii) entraîne i). Il suffit de montrer que $p$ est couvrant
pour la topologie fppf ; or $p$ est fidèlement plat par hypothèse, et
localement de présentation finie, puisque $X$ et $Y$ sont localement de
présentation finie sur $S$.

\textbf{10.9.} — Dans la situation rappelée au début de (10.0), soient
$j \in I$, et $G_{j}$ un $S_{j}$-groupe de présentation finie. Pour que
$G^{\circ} = (G_{j} \times_{S_{j}} S)^{\circ}$ soit représentable, il
faut et il suffit qu'il existe \struck{\ill{}} $i \in I$, $i \geqslant j$,
tel que $(G_{i})^{\circ} = (G_{j} \times_{S_{j}} S_{i})^{\circ}$ soit
représentable.

La condition est suffisante (3.3).

Montrons que la condition est nécessaire. \struck{Alors} \add{En effet,}
$G^{\circ}$ est de présentation finie sur $S$ et ouvert dans $G$ (3.9).
D'après (10.0), il existe $i \in I$, $i \geqslant j$, et un
sous-préschéma en groupes ouvert $G_{i}^{\circ}$ de $G_{i}$ tel que
$G_{i}^{\circ} \times_{S_{i}} S = G^{\circ}$. Le morphisme structural
$G^{\circ} \to S$ est connexe (i.e. à fibres géométriquement connexes)
(Exp. VI$_{\mathrm{A}}$ 2.1.1). Alors, (EGA IV 9.3.3 et 9.7.7), quitte à
augmenter $i$, on peut supposer que le morphisme structural
$G_{i}^{\circ} \to S$ est connexe, donc que l'espace sous-jacent à
$G_{i}^{\circ}$ n'est autre que $(G_{i})^{\circ}$ (3.10.1), donc que
$G_{i}^{\circ}$ représente $(G_{i})^{\circ}$.

\textbf{10.10.} Rappelons deux cas particuliers très utiles de la
situation énoncée en (10.0).
\begin{itemize}
  \item[a)] Etant donné un point $x$ d'un préschéma $X$, on pose
    $S_{0} = \mathrm{Spec}\, \mathbb{Z}$ et on considère la famille
    $(S_{i})_{i \in I}$ filtrante décroissante des voisinages
    \struck{\ill{}} ouverts affines de $x$ ; alors
    $S = \mathrm{Spec}\, \mathcal{O}_{X,x}$. En particulier, si $x$ est le
    point générique d'un schéma intègre $X$, on trouve $S = \kappa(x)$.
  \item[b)] On pose $S_{0} = \mathrm{Spec}\, \mathbb{Z}$, et on considère
    la famille $(\mathcal{A}_{i})_{i \in I}$ préordonnée par inclusion des
    sous-$\mathbb{Z}$-algèbres de type fini de l'anneau d'un schéma affine
    $S$.
\end{itemize}
\page{37}
Etant donné que les $\mathcal{A}_{i}$ sont des anneaux \struck{\ill{}}
noethériens, cela permet dans de nombreux \add{cas} de passer du cas
général au cas noethérien.
\note{« général » et « noethérien » sont cerclés et reliés par une double
flèche, avec dans la marge, cerclé, « intervertir » : lire \emph{du cas
noethérien au cas général}.}

Nous allons maintenant donner deux résultats concernant le cas
particulier envisagé en a).

\textbf{Proposition 10.11.} — Soient $S$ un préschéma intègre de point
générique $\eta$, $G$ un $S$-groupe de présentation finie, $X$ un
$S$-préschéma de présentation finie, $u$ et $v$ deux $S$-morphismes de $X$
dans $G$, $i \colon H \to G$ et $j \colon K \to G$ deux monomorphismes de
présentation finie. Alors :
\begin{itemize}
  \item[(i)] il existe un ouvert non vide $S'$ de $S$ tel que si on note
    $G'$ (resp. $H'$, resp. $X'$, resp. $u'$, resp. $v'$) la restriction
    de $G$ (resp. $H$, resp. $X$, resp. $u$, resp. $v$) au-dessus de
    l'ouvert $S'$ de $S$, le foncteur $\underline{\mathrm{Transp}}(u', v')$
    (resp. $\underline{\mathrm{Centr}}(u')$,
    $\underline{\mathrm{Centr}}_{G'}H'$) soit représentable par un
    sous-$S'$-préschéma fermé de présentation finie $T'_{1}$ (resp.
    $C'_{1}$, resp. $C'_{2}$) de $G'$.
  \item[(ii)] supposons que $K_{\eta}$ soit fermé dans $G_{\eta}$ (ce qui
    est le cas si $K$ est un sous-\struck{$S$}-préschéma en groupes de $G$
    (1.4.2)) ; alors il existe un ouvert non vide $S'$ de $S$ tel que si
    on note $G'$ (resp. $H'$, resp. $K'$) la restriction de $G$ (resp.
    $H$, resp. $K$) au-dessus de l'ouvert $S'$ de $S$, le foncteur
    $\underline{\mathrm{Transp}}_{G'}(H', K')$ (resp.
    $\underline{\mathrm{Norm}}_{G'}K'$) soit représentable par un
    sous-$S'$-préschéma fermé de présentation finie $T'_{2}$ (resp. $N'$)
    de $G'$.
  \item[(iii)] supposons que $H_{\eta}$ et $K_{\eta}$ soient fermés dans
    $G_{\eta}$ (ce qui est le cas si $H$ et $K$ sont des sous-préschémas
    en groupes de $G$ (1.4.2)) ; alors il existe \add{un ouvert} non vide
    $S'$ de $S$ tel que si on note $G'$ (resp. $H'$, resp. $K'$) la
    restriction de $G$ (resp. $H$, resp. $K$) au-dessus de l'ouvert $S'$
    de $S$, le foncteur $\underline{\mathrm{Transpstr}}_{G'}(H', K')$ soit
    représentable par un sous-$S'$-préschéma fermé de présentation finie
    de $G'$.
\end{itemize}

Puisque $G_{\eta}$, $X_{\eta}$, $H_{\eta}$, $K_{\eta}$ sont plats sur le
corps $\kappa(\eta)$, et que
\page{38}
$G_{\eta}$ est séparé sur $\kappa(\eta)$ (Exp. VI$_{\mathrm{A}}$ 0.2),
d'après (10.2, 10.10 a), EGA IV 8.8.2 et 8.10.5), il existe un ouvert
\add{affine} $S'$ de $S$ tel que $G'$ soit un $S'$-groupe de \struck{type}
\add{présentation} fini\add{e}, plat et séparé sur $S'$, et que $X'$, $H'$
et $K'$ soient plats et de présentation finie sur $S'$. Si on suppose
$H_{\eta}$ (resp. $K_{\eta}$) fermé dans $G_{\eta}$, on peut choisir $S'$
de sorte que $H'$ (resp. $K'$) soit fermé dans $G'$. D'après (EGA IV
11.3.15), quitte à remplacer $S'$ par un ouvert affine de $S'$, on peut
supposer que $G'$, $H'$, $K'$ et $X'$ sont essentiellement libres sur $S'$
(au sens de Exp. VIII (6.1)). Il résulte alors de (Exp. VIII 5.5 b) et
e)$^{(*)}$) que, sous les hypothèses de l'énoncé, les foncteurs considérés
sont représentables par des sous-préschémas fermés (donc de présentation
finie sur $S'$) de $G'$.
\marginal{(*) Voir note page 51}
\note{Appel de note et renvoi ajoutés de sa main, le renvoi au bas de la
page, suivi d'un chiffre cerclé, \uncertain{40}. Les « Exp. » de cette
page sont biffés d'un trait à chaque occurrence (« Exp. » → « VI$_A$ »,
« VIII ») ; on les a conservés.}

\textbf{Proposition 10.12.} — Soient $S$ un préschéma intègre, $G$ un
\struck{\ill{}} $S$-groupe de présentation finie, $A$ et $B$ deux
sous-préschémas en groupes de présentation finie de $G$, à fibre
générique lisse, tels que \struck{\ill{}} \add{$A$ soit à fibre générique
connexe} (resp. \struck{\ill{}} $A$ et $B$ soient invariants dans $G$).
Alors il existe un ouvert non vide $S'$ de $S$ et un sous-préschéma en
groupes \add{fermé, $D'$} de présentation finie, à fibres lisses
\struck{et fermé $D'$} de $G' = G|S'$, tel que $D'$ soit \struck{\ill{}} à
fibres connexes (resp. soit invariant dans $G'$), et que $D'$ représente
le faisceau associé, pour la topologie fppf ou fpqc, au préfaisceau en
groupes des commutateurs de $A$ et $B$ dans $G$. En particulier, pour tout
$s \in S'$, on a $D'_{s} = (A_{s}, B_{s})$ avec les notations de (7.2
(vii)).

Soit $\eta$ le point générique de $S$ ; posons
$D'_{\eta} = (A_{\eta}, B_{\eta})$. D'après (7.8) (resp. 7.3 v)),
$D'_{\eta}$ est connexe (resp. invariant dans $G_{\eta}$). D'après (10.1
et 10.10 a)), il existe un ouvert non vide $S'$ de $S$ et un
sous-$S'$-préschéma en groupes $D'$ de présentation finie et fermé dans
$G'$, à fibres séparables (EGA IV 9.7.7 et 9.3.3), ayant $D'_{\eta}$ pour
fibre au point $\eta$, et tel que $D'$ soit à fibres connexes (i.e.
géométriquement connexes (Exp. VI$_{\mathrm{A}}$ 2.1.1)) (resp. soit
invariant dans $G'$) (EGA IV 9.7.7 et 9.3.3 (resp. (10.4)). De plus, nous
avons vu, au cours de la démonstration de (7.8), qu'il existe un entier
$n$ tel que
$\psi_{\eta}^{n}\bigl((A_{\eta} \times_{\kappa(\eta)} B_{\eta})^{n}\bigr)
= D'$, où
$\psi_{\eta} \colon A_{\eta} \times_{\kappa(\eta)} B_{\eta} \to G_{\eta}$
\page{39}
est défini comme en (7.2 (vii)). Nous pouvons définir, par les mêmes
formules qu'en (7.2 (vii) et 7.1 (ii)), les morphismes
$\psi' \colon A' \times_{S'} B' \to G'$ et
$\psi'^{n} \colon (A' \times_{S'} B')^{n} \to G'$, on a bien
$\psi'^{n}(\kappa(\eta)) = \psi_{\eta}^{n}$. Par conséquent (EGA IV 8.8.3
et 8.10.5 \add{et 17.2.6}), on peut choisir $S'$ tel que le morphisme
$\psi'^{n}$ \add{soit plat et} se factorise à travers $D'$, et que le
morphisme $(A' \times_{S'} B')^{n} \to D'$ déduit de $\psi'^{n}$ soit
surjectif. Alors (7.5) le morphisme $(A' \times_{S'} B')^{n+1} \to D'$
déduit de $\psi'^{n+1}$ est couvrant pour la topologie fppf (et a fortiori
pour la topologie fpqc). Donc $D'$ représente le faisceau associé, pour la
topologie fppf ou fpqc, au préfaisceau des commutateurs de $A$ et $B$ dans
$G$. De plus, pour tout $s \in S'$, le morphisme
$(A_{s} \times_{\kappa(s)} B_{s})^{n} \to D'_{s}$ déduit de
$\psi_{s}^{n}$ est surjectif ; donc (7.6), $(A_{s}, B_{s})$ représente le
faisceau des commutateurs de $A_{s}$ et $B_{s}$ dans $G_{s}$ (pour la
topologie fppf ou fpqc), d'où $D'_{s} = (A_{s}, B_{s})$.

\textbf{Corollaire 10.13.} — Soient $S$ un préschéma intègre de point
générique $\eta$, $G$ un $S$-groupe de présentation finie, $D$ un
sous-$S$-préschéma en groupes de présentation finie de $G$ à fibres
lisses et invariant dans $G$. Supposons qu'on ait
$(D_{\eta}, D_{\eta}) = D_{\eta}$ (resp. $(G_{\eta}, D_{\eta}) = D_{\eta}$).
Il existe alors un ouvert non vide $S'$ de $S$ tel que pour tout
$s \in S'$, on ait $(D_{s}, D_{s}) = D_{s}$ (resp.
$(G_{s}, D_{s}) = D_{s}$).

Cela résulte immédiatement de (10.12) et de (EGA IV 8.8.2.5).

Les énoncés (10.\struck{1}\add{2}) et (10.\struck{2}\add{3}) concernant la
catégorie des $S$-groupes de présentation finie s'étendent à la catégorie
des couples formés d'un $S$-groupe de présentation finie et d'un
\struck{\ill{}} $S$-préschéma de présentation finie à groupe d'opérateurs
$G$. De façon précise :

\textbf{10.14.} (i) Dans la situation rappelée au début de (10.0), soient
$j \in I$, $G_{j}$ et $G'_{j}$ deux $S_{j}$-groupes de présentation finie,
$H_{j}$ (resp. $H'_{j}$) un $S_{j}$-préschéma de présentation
\page{40}
finie à groupe d'opérateurs $G_{j}$ (resp. $G'_{j}$). Posons, pour
$i \in I$, $i \geqslant j$, $G_{i} = G_{j} \times_{S_{j}} S_{i}$ et
$G = G_{j} \times_{S_{j}} S$, et définissons de même $G'_{i}$, $G'$,
$H_{i}$, $H$, $H'_{i}$, et $H'$. Notons
$\mathrm{Dihom}_{S\text{-gr.}}((G, H), (G', H'))$ l'ensemble des
di-morphismes de $S$-groupes et de $S$-préschémas à groupe d'opérateurs
du couple $(G, H)$ dans le couple $(G', H')$. Alors l'application
canonique de
$\varinjlim_{i \geqslant j} \mathrm{Dihom}_{S_{i}\text{-gr.}}((G_{i}, H_{i}),
(G'_{i}, H'_{i}))$ dans $\mathrm{Dihom}_{S\text{-gr.}}((G, H), (G', H'))$
est bijective.

(ii) Dans la situation rappelée au début de (10.0), soient $G$ un
$S$-groupe \struck{\ill{}} de présentation finie et $H$ un $S$-préschéma
de présentation finie à groupe d'opérateurs $G$ ; il existe alors un
indice $j \in J$, un $S_{j}$-groupe de présentation finie $G_{j}$, un
$S_{j}$-préschéma de présentation finie $H_{j}$ à groupe d'opérateurs
$G_{j}$ et un di-isomorphisme de $S$-groupes et de $S$-préschémas à
groupes d'opérateurs de $(G_{j} \times_{S_{j}} S, H_{j} \times_{S_{j}} S)$
sur $(G, H)$.
\note{« $j \in J$ » est tapé ainsi, pour $j \in I$ ; tel quel.}

\struck{Proposition} \add{Définition} \textbf{10.15.} — Etant donné un
$S$-foncteur en groupes $G$ et un $S$-foncteur $H$ à groupe d'opérateurs
$G$, on dit que $H$ est un \emph{$S$-foncteur en espaces homogènes sous
$G$ pour la topologie $\mathbb{T}$}, si le morphisme canonique
$G \times_{S} H \to H \times_{S} H$ (défini par
$(g, h) \mapsto (g.h, h)$ pour $g \in G(T)$, $h \in H(T)$) est un
épimorphisme de faisceaux pour la topologie $\mathbb{T}$.

Si $G$ est un $S$-groupe et si $H$ est un $S$-préschéma à
\struck{\ill{}} groupe d'opérateurs $G$, dire que $H$ est un
$S$-préschéma en espaces homogènes sous $G$ pour la topologie
$\mathbb{T}$ revient donc à dire que le morphisme canonique
$G \times_{S} H \to H \times_{S} H$ défini comme précédemment est couvrant
pour la topologie $\mathbb{T}$ (Exp. IV 4.4.3).

De même, dire qu'un $S$-préschéma $H$ à groupe d'opérateurs un $S$-groupe
$G$ est un fibré principal homogène sous $G$ pour la topologie
$\mathbb{T}$ (Exp. IV 5.1.5) revient à dire (Exp. IV 5.6 (ii)) que le
morphisme canonique $G \times_{S} H \to H \times_{S} H$ est un
isomorphisme et que le morphisme structural $H \to S$ est couvrant pour
la topologie $\mathbb{T}$.

\textbf{10.16.} — Dans la situation rappelée au début de (10.0), soient
$j \in I$, $G_{j}$ un $S_{j}$-grou-
\note{Le feuillet s'arrête ici, au milieu du mot ; la suite est dans le
lot suivant.}

\end{document}
